% ==================================================================
\section{Chapter 4 --- Continuous Semimartingales}
% ==================================================================
\setcounter{subsection}{-1}
\subsection{Finite variation calculus (Lebesgue--Stieltjes toolkit)}
\emph{Not Le Gall (he sketches this in passing, \LG{p.\ 71}). This is the lecturer's own build-up (weeks 3.1--3.2), taking finiteness of variation as the starting point rather than deriving it from an alternative construction. Throughout, $a,\tilde a:[0,\infty)\to\R$ are real-valued functions and $\pi=\{0=t_0<t_1<\dots<t_p\}$ ranges over finite partitions with mesh $\|\pi\|=\max_i|t_i-t_{i-1}|$.}

\begin{definition}[Variation, positive/negative variation] \label{chap4.0:definition-variation} \LN{Def 7.1} \\
The \textit{variation function} of $a(\cdot)$ is
$$|a|(t)=\sup_\pi \sum_{i=1}^p\big|a(t_i\wedge t)-a(t_{i-1}\wedge t)\big|, \qquad t\ge0.$$
We say $a(\cdot)$ is of \textit{finite variation} (FV) if $|a|(t)<\infty$ for every $t>0$. The \textit{positive} and \textit{negative variation} are, with $x^+=\max(x,0)$, $x^-=\max(-x,0)$,
$$|a|_+(t)=\sup_\pi\sum_{i=1}^p\big(a(t_i\wedge t)-a(t_{i-1}\wedge t)\big)^+, \qquad |a|_-(t)=\sup_\pi\sum_{i=1}^p\big(a(t_i\wedge t)-a(t_{i-1}\wedge t)\big)^-.$$
Link to Le Gall's notation: $|a|(t)=\int_0^t|da(s)|$ \LG{p.\ 71}.
\end{definition}

\begin{proposition}[Basic properties] \LN{Prop 7.2, 7.3, 7.4} \\
Every monotone function is FV on compacts, with $|a|(t)=|a(t)-a(0)|$. The collection of FV functions is a linear space, and if $a(\cdot)$ is FV then $|a|(\cdot)$, $|a|_+(\cdot)$, $|a|_-(\cdot)$ are all non-decreasing.
\end{proposition}

\begin{theorem}[Jordan decomposition] \label{chap4.0:theorem-jordan_decomposition} \LN{Thm 7.5} \\
$a(\cdot)$ is FV if and only if it can be written as the difference of two non-decreasing functions. The canonical such decomposition is
$$a(t)=a(0)+|a|_+(t)-|a|_-(t), \qquad |a|(t)=|a|_+(t)+|a|_-(t).$$
\end{theorem}

\begin{definition}[CADLAG function] \LN{Def 6.1} \\
$a:[0,\infty)\to\R$ is \textit{CADLAG} (\textit{Continue \`A Droite, Limites \`A Gauche}) if it is right-continuous, $a(t)=a(t+)$, and has finite left limits $a(t-)=\lim_{s\uparrow t}a(s)$ for every $t>0$. Write $\Delta a(t)=a(t)-a(t-)$ for the jump at $t$. CADLAG functions form a linear space, and a CADLAG function has at most countably many jump times \LN{Prop 6.2, 6.3}.
\end{definition}

\begin{proposition}[CADLAG functions of finite variation] \LN{Prop 7.7, 7.9} \\
The CADLAG FV functions form a linear space, and the variation functions $|a|,|a|_\pm$ inherit the CADLAG property, with jumps determined directly: $\Delta|a|(t)=|\Delta a(t)|$ and $\Delta|a|_\pm(t)=(\Delta a(t))^\pm$.
\end{proposition}

\begin{theorem}[Characterisation of FV via partition sums] \label{chap4.0:theorem-fv_characterisation} \LN{Thm 7.8} \\
Let $a(\cdot)$ be CADLAG and $V^a(t;\pi)=\sum_{i=1}^p|a(t_i\wedge t)-a(t_{i-1}\wedge t)|$. For any sequence of partitions $\pi_n$ with $\|\pi_n\|\to0$ and $t^n_{p_n}\to\infty$:
\begin{itemize}
    \item If $a(\cdot)$ is FV, then $V^a(\cdot;\pi_n)\to|a|(\cdot)$ uniformly on $[0,t]$ for every $t\ge0$.
    \item Conversely, if $V^a(\cdot;\pi_n)\to g(\cdot)$ uniformly on $[0,t]$ (for every $t$) for some non-decreasing $g$, then $a(\cdot)$ is FV and $|a|=g$.
\end{itemize}
\end{theorem}

\begin{definition}[Lebesgue--Stieltjes measure and integral] \label{chap4.0:definition-lebesgue_stieltjes} \LN{Def 8.3, 8.7, Prop 8.2, 8.5} \\
For $a(\cdot)$ CADLAG and FV, its Jordan decomposition $a=a(0)+|a|_+-|a|_-$ induces a \textit{signed measure} $\lambda^a$ via
$$\lambda^a((s,t])=\lambda^{|a|_+}((s,t])-\lambda^{|a|_-}((s,t])=a(t)-a(s), \qquad 0\le s<t,$$
where each $\lambda^{|a|_\pm}$ is the (positive) Lebesgue--Stieltjes measure of the non-decreasing $|a|_\pm$. For Borel $\varphi:(0,\infty)\to\R$, the \textit{Lebesgue--Stieltjes integral} is
$$\int_0^t\varphi(s)\,da(s)=\int_0^t\varphi(s)\,d|a|_+(s)-\int_0^t\varphi(s)\,d|a|_-(s),$$
well-defined (``$\varphi$ is $a$-integrable'') provided $\int_0^t|\varphi(s)|\,d|a|(s)<\infty$. It is bilinear in $(\varphi,a)$, and $\lambda^a(\{t\})=\Delta a(t)$, $\int_0^t da(u)=a(t)-a(0)$.
\end{definition}

\begin{proposition}[$a$-absolute continuity] \label{chap4.0:proposition-a_absolute_continuity} \LN{Def 8.10, Prop 8.11} \\
Let $a(\cdot)$ be CADLAG FV and $\varphi(\cdot)$ be $a$-integrable, and define $y(t)=y_0+\int_0^t\varphi(s)\,da(s)$. Then $y(\cdot)$ is CADLAG FV, $\Delta y(t)=\varphi(t)\Delta a(t)$, and $y(\cdot)$ is continuous whenever $a(\cdot)$ is. Chain rule: if $\psi(\cdot)$ is $y$-integrable, $\int_0^t\psi\,dy=\int_0^t\psi\varphi\,da$.
\end{proposition}

\begin{theorem}[Integration by parts / Itô formula for FV functions] \label{chap4.0:theorem-ibp_ito_fv} \LN{Thm 9.1, 9.2} \\
Let $a,\tilde a$ be CADLAG FV and $f\in C^1(\R)$. Then
$$a(t)\tilde a(t)-a(0)\tilde a(0)=\int_0^t a(s)\,d\tilde a(s)+\int_0^t\tilde a(s-)\,da(s),$$
$$f(a(t))-f(a(0))=\int_0^t f'(a(s-))\,da(s)+\sum_{0<s\le t}\Big(f(a(s))-f(a(s-))-f'(a(s-))\Delta a(s)\Big).$$
\end{theorem}
\begin{proofidea} Approximate by Riemann--Stieltjes sums along a telescoping partition and pass to the limit; the correction sum collects exactly the second-order error at each jump, which vanishes automatically if $a(\cdot)$ is continuous. \end{proofidea}

\begin{trap} This pathwise Lebesgue--Stieltjes machinery \emph{requires} the integrator to be FV: if the Riemann--Stieltjes sums $\sum\varphi(t^n_{i-1})(a(t^n_i\wedge t)-a(t^n_{i-1}\wedge t))$ converge to the same limit along \emph{every} continuous $\varphi$ and every mesh$\to0$ partition sequence, then $a(\cdot)$ must be FV \LG{Lemma 4.3}. Since Brownian paths are a.s.\ not FV (\cref{sec:qv} recap below), this rules out defining $\int H\,dB$ pathwise --- this is precisely the gap Itô calculus (Ch.\ 5) fills. \end{trap}
\begin{quizbox} Jordan decomposition statement (\cref{chap4.0:theorem-jordan_decomposition}); that CADLAG FV functions form a linear space; the FV Itô formula (\cref{chap4.0:theorem-ibp_ito_fv}) with its jump-correction sum --- this is the deterministic skeleton the stochastic Itô formula in Ch.\ 5 generalises, and for continuous $a(\cdot)$ it collapses to ordinary integration by parts with \emph{no} bracket term (unlike Itô's formula for BM). \end{quizbox}

\subsection{Finite variation processes}
\begin{definition}[Continuous process of finite variation] \label{chap4.1:definition-continuous_fv_process} \LG{Def 4.4} \\
A process $A=(A_t,t\ge0)$ is a \textit{continuous process of finite variation} if $A_0=0$ a.s., $A$ is adapted to $(\mc F_t)$, and (a.s.) every sample path $t\mapsto A_t(\omega)$ is continuous and of finite variation. It is a \textit{continuous increasing process} if every sample path is non-decreasing.
\end{definition}

\begin{proposition}[Integral process] \label{chap4.1:proposition-integral_process} \LG{Prop 4.5} \\
Let $A$ be a continuous process of finite variation and $H=(H_t,t\ge0)$ progressively measurable. Define the integral process path-by-path,
$$(H\cdot A)_t(\omega)=\int_0^t H_s(\omega)\,dA_s(\omega), \qquad t\ge0,$$
well-defined whenever $\int_0^t|H_s(\omega)|\,d|A|_s(\omega)<\infty$ for every $t$ and (a.e.) $\omega$. Then $H\cdot A$ is again a continuous process of finite variation, and the chain rule holds: $K\cdot(H\cdot A)=(KH)\cdot A$ whenever $K$ is integrable w.r.t.\ $H\cdot A$.
\end{proposition}

\begin{definition}[Localizing sequence, local consistency] \label{chap4.1:definition-localizing_sequence} \LN{p.\ 76} \\
A \textit{localizing sequence} is a non-decreasing sequence of stopping times $(T_n,n\in\N)$ with $T_n(\omega)\uparrow\infty$ for every $\omega$. If a sequence of processes $(X^{(n)})$ is \textit{consistent under stopping} along $(T_n)$, i.e.\ $X^{(m)}_{T_m\wedge t}=X^{(n)}_{T_m\wedge t}$ for $m\le n$ and all $t$, then there is a unique process $X$ such that $X_t=\lim_n X_t^{(n)}$ for all $t$ and $X_{T_n\wedge t}=X^{(n)}_{T_n\wedge t}$ for every $n,t$; $X$ inherits any property (adaptedness, continuity, monotonicity) shared by every $X^{(n)}$.
\end{definition}
\begin{usedin} This \textit{local consistency lemma} (together with the deterministic version for a single sequence of stopping times / paths) is the standard gluing device: it is reused verbatim to construct the quadratic variation of a continuous local martingale (\cref{sec:qv} below, via $T_n=\inf\{t:|M_t|\ge n\}$) and to extend the stochastic integral from $\HH$ to $L^2_{\mathrm{loc}}(M)$ in Ch.\ 5. \end{usedin}
\begin{quizbox} Stable-under-stopping list: adapted + right-continuous, continuous FV process, right-continuous martingale, and UI right-continuous martingale are all preserved under $X\mapsto X^T$ (\cref{chap4.2:proposition-stable_under_stopping} below extends this to local martingales). Know this list --- ``is $X^T$ still in the same class?'' is a recurring quiz question. \end{quizbox}

\subsection{Continuous local martingales}
Once again we consider a filtered probability space $(\Omega, \mc{F}, (\mc{F}_t), P)$.

\begin{definition}[Stopped process] \label{chap4.2:definition-stopped_process}
\LG{p. 75, start of chapter 4.2} \\
Let $T$ be a stopping time, and let $X = (X_t)_{t\geq0}$ be an adapted process with continuous sample paths. The \textit{process $X$ stopped at $T$} is the process $X^T=(X_t^T)_{t \geq 0}$ defined by 
$$X_t^T=X_{t \land T}$$
for every $t \geq 0$.
\end{definition}

\begin{proposition}\label{chap4.2:proposition-stable_under_stopping}
\LG{p. 75, start of chapter 4.2} \\
Let $S$ and $T$ be stopping times, and let $X = (X_t)_{t\geq0}$ be an adapted process with continuous sample paths. Then
$$(X^T)^S=(X^S)^T=X^{S \land T}$$
\end{proposition}
\begin{proofidea}
Use the commutative properties of $\land$.
\end{proofidea}


\begin{definition} [Continuous local martingale] \label{chap4.2:definition-continuous-local-martingale}
\LG{Definition 4.6} \\
Let $M=(M_t)_{t \geq 0}$ be a random process with continuous sample paths with $M_0=0$ a.s. If there exists a non-decreasing countable sequence $(T_n)_{n \in \N_0}$ of stopping times such that for every $\omega \in \Omega$ then $T_n(\omega) \uparrow\infty$ (written $T_n \uparrow\infty$) and, for every $n$, the stopped process $M^{T_n}$ is uniformly integrable (note property (b) below!), then we call $M$ a \textit{continuous local martingale}.

More generally, when we do not assume that $M_0=0$ a.s., we say that $M$ is a continous local martingale if the process $N_t=M_t-M_0$ is a continous local martingale. 

In all cases, we say that the sequence of stopping times $(T_n)$ \textit{reduces} $M$ if $T_n \uparrow \infty$ and, for every $n$, the stopped process $M^{T_n}$ is a uniformly integrable martingale.
\end{definition}

\begin{remark}[Remarks on chosen definition] \label{chap4.2:remark-remarks_on_definition_of_continous_local_martingale}
\LG{p. 76}
\begin{enumerate}[(i)]
    \item We do not require that $M_t$ are in $L^1$ in the definition of a continuous local martingale. In particular, the variable $M_0$ may be any $\mc{F}_0$-measurable random variable, even though the continuity of the process disallows very degenerate distributions (for example it disallows $M_0=\infty$ a.s., because then $M$ would not be continuous).
    \item Any martingale with continuous sample paths is a continuous local martingale but the converse is false, and for this reason we will sometimes speak of \textit{true martingales} to emphasize the difference with local martingales. Examples:
    \begin{enumerate}[label=(Ex.\arabic*.)]
        \item Let $B$ be an ($\mc{F}_t$)-Brownian motion started from 0, and let $Z$ be an $\mc{F}_0$-measurable random variables. Then the process $M_t = Z + B_t$ is always a continuous local martingale, but it is not a martingale if $\E[|Z|]=\infty$
        \item If we require $M_0=0$, we can consider $M_t=ZB_t$, which is always a continuous local martingale (Exercise 4.22 in Le Gall or 5.3 (a) in problem sheets for the course) but is not a martingale if $\E[|Z|]=\infty$.
    \end{enumerate}
    \item One can define a notion of local martingale with càdlàg sample paths. In this course, however, we consider only continuous local martingales.
\end{enumerate}
\end{remark}
\begin{trap} Local martingale $\not\Rightarrow$ martingale, even UI (e.g.\ inverse Bessel-3). \end{trap}

\begin{properties}
\LG{pp. 76-77}
\begin{enumerate}[(a)]
    \item A (true) martingale with continous sample paths is a continous local martingale, and the sequende $T_n=n$ reduces $M$.
    \item In the definition of a continous local martingale starting from 0, one can replace "uniformly integrable martingale" by merely "martingale".
    \item If $M$ is a continous local martingale, then, for every stopping time $T$, then $M^T$ is a continouous local martingale. (This follows from Corollary 3.24 (\todo{Add ref})).
    \item If $(T_n)$ reduces $M$ and if $(S_n)$ is a sequence of stopping times such that $S_n \uparrow \infty$, then the sequence $(T_n \land S_n)$ also reduces $M$ (This follows from Corollary 3.24 again (\todo{Add ref})).
\end{enumerate}
\end{properties}
\begin{proof}
\todo{add} 
\end{proof}


\begin{proposition}[Three useful properties] 
\LG{pp. 77-78}
\label{chap4.2:proposition-three_useful_properties_of_continous_local_martingales}
\begin{enumerate}[(i)]
    \item A non-negative continous local martingale $M$ such that $M_0 \in L^1$ is a supermartingale.
    \item A continuous local martingale $M$ such that there exists a random variable $Z \in L^1$ with $|M_t| \leq Z$ for every $t \geq 0$ (in particular a bounded continuous local martingale) is a uniformly integrable martingale.
    \item If $M$ is a continous local martingale and $M_0=0$ ( or more generally $M_0 \in L^1$), the sequence of stopping times
    $$T_n = \inf\{t \geq 0 | \;|M_t| \geq n\}$$
    reduces M.
\end{enumerate}
\end{proposition}
\begin{proofidea}
Write $M_t=M_0+N_t$. The idea is to use the sequence $(T_n)$ of stopping times that reduces $N_t$ to the martingale property, adding $M_0$ on both sides to see
$$M_{s \land T_n}=\E[M_{t \land T_n}| \mc{F}_s], \quad s \leq t,$$
Using Fatou's Lemma in (i) shows the supermartingale property. Using dominated convergence shows in (ii) shows $L^1$ convergence $M_{t \land T_n} \to M_t$, which can be used to show the martingale property. 
For (iii) then observe that $M^{T_n}$ is dominated by $n+|M_0|$ and use (ii).
\end{proofidea}
\begin{trap}
Considering property (ii) in Proposition \ref{chap4.2:proposition-three_useful_properties_of_continous_local_martingales}, one might expect that a continuous local martingale $M$ such that the collection $(M_t)_{t\geq 0}$ is uniformly integrable is automatically a martingale. This is incorrect! This condition is slightly weaker than in the Proposition, and, for instance, if $B$ is a three-dimensional Brownian Motion started from $x\neq 0$, then the process $M_t=1/|B_t|$ is a continous local martingale bounded in $L^2$, but it is not a martingale. \\
In general not even the slightly stronger condition a continuous local martingale satisfies the stronger property of being bounded in $L^p$ for some $p>1$ can save the statement in the generality.
\end{trap}

Now we observe that the classes of finite variation processes and continuous local martingales are distinct classes with one unique (degenerate) common member: 
\begin{theorem}\label{chap4.2:theorem-continous_local_martingale_plus_finite_variation_process_implies_zero_process}
\LG{Theorem 4.8}
Let $M$ be a continous local martingale. Assume that $M$ is also a finite variation process with $M_0=0$. Then $M_t = 0$ almost surely for every $t \geq 0$.
\end{theorem}
\begin{quizbox} This theorem is the uniqueness engine for $\br{M}$ and the semimartingale decomposition. \end{quizbox}
\begin{proofidea}
Set $\tau_n = \inf\{t \geq 0| \; \int_0^t|dM_s| \geq n\}$
for every integer $n \in \N_0$. This stopping time can for a fixed $n$ show that for the process $N=M^{\tau_n}$, that
$$\E[N^2_t] \leq n \E\left[\sup_{1 \leq i \leq p} |N_{t_i}-N_{t_{i-1}}| \right]$$
for some subdivisions. Using a sequence of subdivisions whose mesh tends to 0, shows that $\E[N^2_t]=0$ and hence $M_{t \land \tau_n}=0$ a.s. Letting $n \to \infty$ shows $M_t=0$ a.s.
\end{proofidea}

\begin{corollary}\label{chap4.2:corollary-continous_FV_local_martingale_is_constant}
Let $M$ be a continuous local martingale (not necessarily with $M_0=0$) which is also a finite variation process. Then $M_t=M_0$ almost surely, for every $t \geq 0$; that is, $M$ is a.s.\ constant.
\end{corollary}

\begin{proofidea}
By definition, $N_t:=M_t-M_0$ is a continuous local martingale with $N_0=0$. Since $t \mapsto M_0$ is constant (hence FV) and $M$ is FV, $N=M-M_0$ is also a finite variation process (difference of two FV processes, with the total variation of $N$ on $[0,t]$ equal to that of $M$). Theorem \ref{chap4.2:theorem-continous_local_martingale_plus_finite_variation_process_implies_zero_process} applied to $N$ gives $N_t=0$ a.s.\ for every $t$, i.e.\ $M_t=M_0$ a.s.
\end{proofidea}
\begin{usedin} This is precisely the ``uniqueness engine'' cited in the toolbox: it gives uniqueness of the semimartingale decomposition $X=X_0+M+A$ (\cref{sec:continous_semimartingales}) and, together with Lévy's characterisation, uniqueness of the quadratic variation $\br{M}$. \end{usedin}



\subsection{Quadratic variation of a continuous local martingale}\label{sec:qv}
\begin{theorem}[Existence of $\br{M}$] \todo{unique increasing continuous $\br{M}$, $\br{M}_0=0$, $M^2-\br{M}$ loc.\ mart; limit of $\sum(M_{t_i}-M_{t_{i-1}})^2$ in probability} \end{theorem}
\begin{proofidea} \todo{bounded case first, $L^2$-Cauchy via Doob; localise} \end{proofidea}
\begin{proposition} \todo{$\br{M^T}=\br{M}^T$; $M$ constant on $[a,b]$ iff $\br{M}$ is} \end{proposition}
\begin{theorem}[$L^2$-bounded martingales] \label{chap4.3:theorem-L2_bounded_martingales}
\LG{Theorem 4.10 --- check numbering against your edition} \\
Let $M$ be a continuous local martingale with $M_0 \in L^2$. Then $M$ is a (uniformly integrable, $L^2$-bounded) martingale of class $\HH$ if and only if $\E[\br{M}_\infty] < \infty$. In that case:
\begin{enumerate}[(i)]
    \item $M_t^2-\br{M}_t$ is a uniformly integrable martingale (not just a local martingale);
    \item $\E[M_\infty^2]=\E[M_0^2]+\E[\br{M}_\infty]$, and more generally $\E[M_t^2]=\E[M_0^2]+\E[\br{M}_t]$ for every $t$.
\end{enumerate}
\end{theorem}
\begin{proofidea}
($\Leftarrow$) Localise with $T_n=\inf\{t:\br{M}_t\ge n\}$; on $M^{T_n}$ the process $(M^{T_n})^2-\br{M}^{T_n}$ is a true martingale (bounded increments in expectation), which gives $\E[(M_{t\wedge T_n})^2]=\E[M_0^2]+\E[\br{M}_{t\wedge T_n}]\le \E[M_0^2]+n$, so $M^{T_n}$ is $L^2$-bounded uniformly in $n$. Fatou plus $\E[\br M_\infty]<\infty$ (dominated convergence along $T_n\uparrow\infty$) upgrades this to $M\in\HH$ and to the identity in (ii); (i) then follows since $M^2-\br M$ is a local martingale dominated appropriately, hence UI by \cref{chap4.2:proposition-three_useful_properties_of_continous_local_martingales}(ii)-type domination.
($\Rightarrow$) If $M\in\HH$, apply the reverse estimate to $M^{T_n}$ and use monotone convergence $\br{M}_{t\wedge T_n}\uparrow\br{M}_t$ together with the $L^2$-bound $\sup_t\E[M_t^2]<\infty$ to conclude $\E[\br{M}_\infty]<\infty$.
\end{proofidea}
\begin{corollary}
If $\E[\br{M}_t]<\infty$ for every $t\ge0$ (not necessarily $\E[\br M_\infty]<\infty$), then $M$ is a true martingale.
\end{corollary}
\begin{proofidea}
Fix $t$ and apply \cref{chap4.3:theorem-L2_bounded_martingales} on $[0,t]$: the stopped process $M^t$ has $\E[\br{M^t}_\infty]=\E[\br{M}_t]<\infty$, so $M^t\in\HH$, in particular $M^t$ is a true (UI) martingale. Since $t$ was arbitrary and $M^t_s=M_s$ for $s\le t$, $M$ itself satisfies the martingale property on every finite horizon.
\end{proofidea}
\begin{definition}[Bracket $\br{M,N}$] \todo{polarisation; bilinear, symmetric; limit of cross-products} \end{definition}
\begin{proposition}[Kunita--Watanabe] \todo{$\int|H||K|\,|d\br{M,N}|\le(\int H^2d\br{M})^{1/2}(\int K^2d\br{N})^{1/2}$} \end{proposition}
\begin{quizbox} $\br{M}$ characterisation; the $\HH$ criterion; KW. \end{quizbox}
\begin{usedin} $\HH$ is the Hilbert space of Ch.\ 5; KW gives $\br{H\cdot M,N}$ well-defined. \end{usedin}

\subsection{Continuous semimartingales}\label{sec:continous_semimartingales}
\begin{definition} \todo{$X=X_0+M+A$; uniqueness of decomposition} \end{definition}
\begin{definition} \todo{$\br{X,Y}:=\br{M,N}$} \end{definition}
\begin{proposition} \todo{$\br{X,Y}$ is limit of $\sum\Delta X\Delta Y$; FV parts contribute nothing} \end{proposition}
